% !TeX spellcheck = en_GB
% Template for Federated Learning Course Project Report

\documentclass[9pt]{article}
\usepackage{spconf,amsmath,bm,graphicx}
\usepackage[a4paper, margin=1in]{geometry}
\usepackage[dvipsnames]{xcolor}
\usepackage{hyperref, cleveref, tcolorbox}
\usepackage{algorithm, algorithmic}
\usepackage{amsfonts,amssymb,amsbsy}
\usepackage{subcaption, adjustbox}
\usepackage{listings}
\usepackage{qrcode}


\newtcolorbox{hintbox}{
	colback=blue!5,
	colframe=blue!60!black,
	boxrule=0.8pt,
	arc=2mm,
	left=1mm,
	right=2mm,
	top=1mm,
	bottom=1mm,
	title=\textbf{Hint (remove before submission)}
}



% Title
\title{Template for a Federated Learning Project Report}
\name{Your Name}
\address{}

\begin{document}
	


	\maketitle
	
	\begin{hintbox}
		All hint boxes in this template are \emph{instructions for drafting} and must be
		\textbf{removed before final submission}. Not every hint corresponds to a
		separate grading criterion; however, ignoring hints typically leads to unclear
		problem formulations, weak experiments, or low review scores.
	\end{hintbox}
	
\begin{abstract}
	\begin{hintbox}
The abstract provides a concise summary of the project, including the FL application, 
empirical graph modeling, variation minimization approach, and the FL algorithms used. 
Do not include detailed numerical results. Instead, focus on problem, method, and evaluation setup.
\end{hintbox}
\end{abstract}
	
\textbf{Keywords:} Federated learning, networks, personalized machine learning, trustworthy artificial intelligence
	
\begin{hintbox}
	Your report must adhere to the following formatting and length requirements:
	\noindent
\begin{itemize}
	\item Do not include any information (GitHub usernames, repository links, or dataset paths) 
	           in the report that could be used to identify you. 
	\item Use the \LaTeX\ settings as defined in this file without changing margins or font sizes.
	\item You must not modify the number or titles of sections in this template. 
	\item The report must not exceed 5 pages in total. The 5\textsuperscript{th} page must 
	contain only the references.
		\end{itemize}
\end{hintbox}

\begin{hintbox}
Students are strongly encouraged to use terminology consistent with the
Aalto Dictionary of Machine Learning \cite{AaltoDictML}.
\end{hintbox}


\section{Introduction}

\begin{hintbox}
	Describe the application scenario, motivation, and learning objective of your 
	federated learning (FL) system.
\end{hintbox}

	
\section{Problem Formulation}
\label{sec:pf}
\begin{hintbox}
	Clearly define the components of your federated learning (FL) network.
	In particular, specify:
	\begin{itemize}
		\item what constitutes a \emph{node} (client) in your system,
		\item how edges between nodes are obtained and how edge weights are defined
		(e.g., geographical proximity, data similarity, or a given communication graph),
		\item See \cite[Ch. 7]{Jung2025} for data-driven methods to choose the edges of an FL network. 
		\item what data each node holds locally (local datasets),
		\item what model is trained locally at each node,
		\item and how the local loss functions are defined.
	\end{itemize}
\end{hintbox}

\begin{hintbox}
	To ensures comparability and fair peer review, all projects must 
	use the same primary data source: the
	\textbf{Finnish Meteorological Institute (FMI) Open Data service}
	(\url{https://en.ilmatieteenlaitos.fi/open-data}). 
	You can find a Python script that demonstrates how to use this data service here: \qrcode{https://github.com/alexjungaalto/FederatedLearning/blob/main/Edition2026/assets/GetFMIData.py}
	
\end{hintbox}

	
\section{Methodology}
\label{sec_meth} 
\begin{hintbox}
The project requires you to apply GTVMin-based methods to 
the FL application modelled in Section \ref{sec:pf}. In this section 
you need to clearly state and explain:
\begin{itemize}
	\item Your choice of variation measure, e.g., $\mathrm{GTV}(\mathbf{w}^{(i)}-\mathbf{w}^{(j)})$
	for parametric models.
	\item Your choice of FL algorithm (i.e., optimization method for solving 
	GTVMin) and its message passing implementation.
\end{itemize}
\end{hintbox}

\begin{hintbox}
You are not expected to propose novel FL algorithms; correct implementation and analysis of existing GTVMin-based methods is sufficient.
\end{hintbox}
	
\section{Numerical Experiments}
\label{sec:experiments}

\begin{hintbox}
	Present quantitative results (training, validation, and test losses for each node of the FL network), 
	comparisons, and figures.
	At a minimum, you must evaluate and compare \textbf{two distinct FL systems}. 
	These systems should differ in a \emph{meaningful design choice}, such as:
	\begin{itemize}
		\item different hyper-parameters of the optimisation method (e.g., regularisation strength, step size),
		\item different graph constructions or edge weights,
		\item different local model classes or model capacities.
	\end{itemize}
	Clearly explain evaluation metrics (e.g., which loss function has been used 
	for training, validation, and final evaluation on the test set) and experimental 
	settings (e.g., how data is split into train/validation/test sets and how hyper-parameters 
	are selected).
\end{hintbox}

\begin{hintbox}
	\textbf{Figures and tables.} 
	Figures should be designed to communicate quantitative results clearly and efficiently.
	Avoid decorative elements that do not convey information (e.g., excessive colours, 3D effects, 
	or chart junk). Axes, legends, and captions must be self-contained and precise. 
	Whenever possible, embed numerical values directly into figures or tables rather 
	than relying on verbose textual explanations. For general principles of effective 
	quantitative graphics, see \cite{TufteVDQI}.
\end{hintbox}



\begin{hintbox}
\textbf{Figures and tables.} 
Figures should be designed to communicate quantitative results clearly and efficiently.
Avoid decorative elements that do not convey information (e.g., excessive colours, 3D effects, 
or chart junk). Axes, legends, and captions must be self-contained and precise. 
Whenever possible, embed numerical values directly into figures or tables rather 
than relying on verbose textual explanations. For general principles of effective 
quantitative graphics, see \cite{TufteVDQI}
\end{hintbox}

\begin{hintbox} 
Your submission must include a zip archive containing a single Python script along with any 
necessary data files. Minimize the use of non-standard Python packages to ensure ease of 
execution and reproducibility.
\end{hintbox} 


	
\section{Conclusion}
	\label{sec:conclusion}
\begin{hintbox}
Interpret the numerical results from Sec. \ref{sec:experiments}. 
In particular, 
\begin{itemize}
	\item Discuss whether the obtained results solve the problem satisfactorily.
	\item Identify limitations of the studied FL methods and suggest potential improvements (e.g., 
	collecting more data points or using larger local models at some nodes).
\end{itemize}
You can find some guidance on how to diagnose ML models in \cite[Sec.~6.6]{Jung2022}.
\end{hintbox}



%	\section*{References}
	\begin{thebibliography}{9}
		\bibitem{Jung2025} A. Jung, \textit{Federated Learning: From Theory to Practice}, Springer, 2026. Arxiv Preprint: \url{https://arxiv.org/abs/2505.19183}.
		\bibitem{Jung2022} A. Jung, \textit{Machine Learning: The Basics}, Springer, 2022. Available via Aalto Library. Preprint: \url{https://github.com/alexjungaalto/MachineLearningTheBasics/}. 
		\bibitem{AaltoDictML} A. Jung (Editor),  \textit{Aalto Dictionary of Machine Learning}, \url{https://aaltodictionaryofml.github.io/}, 2026.
		\bibitem{TufteVDQI}
		E.~R.~Tufte,
		\emph{The Visual Display of Quantitative Information},
		2nd edition, Graphics Press, 2001.
	\end{thebibliography}
	
\end{document}